\chapter{Szeregi formalne i wielomiany.}

\section{Pierścień szeregów formalnych.}
\begin{definition}[Pierścień szeregów formalnych.]
	Niech $R$ będzie pierścieniem przemiennym z $\mathbf{1}$. 
	
	\begin{enumerate}[label=\arabic*)]
		\item 
	Definiujemy $R\left[ \left| X \right|  \right] $jest \textbf{pierścieniem szeregów formalnych } (o współczynnikach z $R$ ) jako zbiór $R^{\N}$, tzn
	\[
	R\left[ \left| X \right|  \right] := R^{\N}
	.\] 

	\noindent \textbf{Dodawanie} jest zdefiniowane jak w $\text{Fun}\left( \N,R \right) $, czyli
	\[
		\left( a_{n} \right)_{n=1}^{\infty} +  \left( b_{n}  \right) _{n=1}^{\infty} := \left( a_{n} +b_{n}  \right) _{n=1}^{\infty}
	.\] 

	\noindent  Mnożenie jest inne, tzw splot:
	\[              
		\left( a_{n} \right)_{n=1}^{\infty} \cdot   \left( b_{n}  \right) _{n=1}^{\infty} := \left( c_{n}  \right) _{n=1}^{\infty}
	,\] 
	gdzie
	\[
		c_{n} = a_0b_{n} + a_1b_{n-1} + \ldots + a_{n}b_{0} = \sum_{i+j = n}^{\infty} a_{i}b_{j},
	.\] 
	\textbf{Oznaczamy:} $\left( a_{n} \right)_{n\in \N}$ oznaczamy przez
	\[
	\sum_{n = 1}^{\infty} a_{n}x^{n}\text{ lub } a_{0}X^{0} + a_{1}X^{1} + \ldots
	.\] 
	To tylko oznaczenie, np. 
	\[
	X = 1\cdot X = \left( 0,1,0,0,\ldots \right) 
	.\] 
	Zadanie:
	\begin{itemize}
		\item $\left( R\left[ \left| X \right|  \right] , +, \cdot  \right) $- pierścień przemienny $\mathbf{1}$,
		\item $R$- dziedzina $\implies R\left[ \left| X \right|  \right] $ też.
	\end{itemize}
	Dowód jako zadanie.
		\item 
			\[
			R\left[ X \right] := \{\sum_{n=0}^{\infty} a_{n} X^{n}\in R\left[ \left| X \right|  \right]: \left(\exists N \right)\left(\forall n \ge N \right) \left( a_{n} = \mathbf{0} \right)  \} 
			.\] 
			(prawie wszystkie zerowe)

			Zadanie: $R\left[ X \right] \underset{\scriptscriptstyle\text{podp}}{\subseteq} R\left[ \left| X \right|  \right].$
			
			Niech $f\in R\left[ X \right] $, 
			\[
			f = \sum_{n=0}^{N} a_{n} X^{n} = a_0 + a_1X + a_2X^2 + \ldots + a_{N}X^{N} + 0\cdot X^{N+1} + 0\cdot X^{N+2} + \ldots
			.\] 
			Każdy wielomian f jak wyżej wyznacza \textbf{funkcję wielomianową}:
			\[
			\overline{f}: R\to R \text{ zadaną przez } \overline{f}\left( r \right) := a_0 + a_1r + a_2r + \ldots + a_{N}r^{N}
			.\] 
	\end{enumerate}
\end{definition} 
\begin{remark}
	.

	\begin{enumerate}[label=\arabic*)]
		\item Różne wielomiany mogą dawać te same funkcje wielomianowe. Na przykład
			\[
			\text{dla } R = \mathbb{Z} _{2}: \overline{X}= \overline{X^2} = \overline{X^3} = \id_{\scriptscriptstyle \mathbb{Z}_{2}}
			.\] 
			bo
			\[
			\begin{cases}
				0 &= 0^2 = 0^3 \\
				1 &= 1^2 = 1^3
			\end{cases}
			.\] 
			Zatem \textbf{nie należy} utożsamiać i mylić wielomianów z funkcjami wielomianowymi
		\item Szereg formalny zwykle nie wyznacza żadnej funkcji, chyba że w $R$ mamy pojęcie zbieżności, na przykład w $R\left[ \left| X \right|  \right] $ jest \textbf{podpierścieniem} szeregów zbieżnych (wszędzie), na przykład 
			\[
			\exp = \sum_{n=0}^{\infty} \frac{1}{n!}x^{n}
			.\] 
	\end{enumerate}
\end{remark} 
\begin{definition}
	Niech $f = \sum_{n=0}^{N} a_{n} X^{n}\in R\left[ X \right] \setminus \{\mathbf{0}\} $.
	\begin{enumerate}[label=\arabic*)]
		\item $\text{deg}\left( f \right):= \max \{n: a_{n} \neq \mathbf{0}\}.$ Wielomian zerowy nie ma stopnia,
		\item $a_0$- \textbf{wyraz wolny} $f$,
		\item $a_{n}$- \textbf{współczynnik wiodący}, gdzie $n = \text{deg}\left( f \right) $
	\end{enumerate}
\end{definition} 
\begin{definition}
	.
	\begin{enumerate}[label=\arabic*)]
		\item $R\left[ \left| X,Y \right|  \right] := R\left[ \left| X \right|  \right] R\left[ \left| Y \right|  \right] $: pierścień szeregów formalnych dwu zmiennych.
		\item $R\left[ \left| X_1,X_2,\ldots, X_{n-1},X_{n}  \right|  \right] := R\left[ \left| X_1,\ldots,X_{n-1} \right|  \right] \left[ \left| X_{n}  \right|  \right] $: pierścień szeregów formalnych $n$ zmiennych.

			$\rotatein$

			 $\sum_{(i_1,\ldots,i_{n})\in \N^{n}} a_{i_1 \ldots i_{n}}X_1^{i_1}\ldots X_{n}^{i_{n}}$: funkcja $f$ z $\N^{n}$ w $R: f\left( i_1,\ldots,i_{n} \right) = a_{i_1\ldots i_{n}}$
		 \item Podobnie $R\left[ X_1,\ldots,X_{n}  \right] = R\left[  X_1,\ldots,X_{n} \right]\left[ X_{n} \right]  $ : pierścień wielomianów $n$-zmiennych.
	\end{enumerate}
\end{definition} 

\section{Pierścienie i ciała liczbowe.}
\begin{definition}[Pierścień Gaussa.]
	\begin{enumerate}[label=\arabic*)]
		\item $\mathbb{Z} \left[ i \right] := \{a + bi: a,b \in \mathbb{Z} \} \subseteq \mathbb{C} $ : \textbf{pierścień Gaussa} 
		\item Niech $\alpha \in \mathbb{C} \setminus \mathbb{Q} $ oraz $\alpha ^2 \in \mathbb{Q} $, definiujemy
			\[
			\mathbb{Q} \left(\alpha  \right) := \{a + b\alpha \in \mathbb{C} : a,b \in \mathbb{Q} \} \underset{\scriptscriptstyle\text{podc}}{\subseteq} \mathbb{C} 
			.\] 
	\end{enumerate}
\end{definition} 
\begin{remark}
	Ogólniej:

	Niech $\alpha \in \mathbb{C} \setminus \mathbb{Z} $, takie że $\alpha ^2 \in \mathbb{Z} $, wtedy
	\[
	\mathbb{Z} \left[ \alpha  \right] := \{a + b\alpha :a,b \in \mathbb{Z} \} \underset{\scriptscriptstyle\text{podc}}{\subseteq} \mathbb{C} \in \mathbb{Z} 
	,\]
	bo:
	\begin{itemize}
		\item podciało (trywialne?)
		\item $\left( a + b\alpha  \right)\left( c + d \alpha  \right) = \left(\underbrace{ac + bd\alpha ^2}_{\in \mathbb{Z} }\right) + \left( \underbrace{ad + bc}_{\in \mathbb{Z} } \right) \alpha \in \mathbb{Z} \left[ \alpha  \right] $
	\end{itemize}
\end{remark} 
\begin{definition}[Homomorfizm pierścieni.]
	Niech $R,S$ będą pierścieniami. Funkcja $f:R \to S$ jest \textbf{homomorfizmem (pierścieni)}, gdy
	\[
	\left(\forall x,y \in R \right)\left( \begin{cases}
			f\left( x +_{\scriptscriptstyle R} y \right) &= f\left( x \right) +_{\scriptscriptstyle S} f\left( y \right) \\
			f\left( x \cdot_{\scriptscriptstyle R} y \right) &= f\left( x \right) \cdot_{\scriptscriptstyle S} f\left( y \right) 
	\end{cases} \right)  
	.\] 

	$f$ jest epimorfizmem, monomorfizmem, izomorfizmem, endomorfizmem, automorfizmem- jak dla grup $R\cong S$, gdy intuicyjnie $f: R \overset{\text{izo}}{\to} S$
\end{definition} 
\begin{remark}
	Przypomnijmy:
	\begin{itemize}
		\item Złożenie homomorfizmów jest homomorfizmem.
		\item Funkcja odwrotna do izomorfizmu jest izomorfizmem.
	\end{itemize}
\end{remark} 
\begin{remark}
	Załóżmy, że $f: R\to S $ jest homomorfizmem pierścieni. Wtedy
	\begin{enumerate}[label=\arabic*)]
		\item $f: \left( R, +_{\scriptscriptstyle R} \right) \to \left( S, +_{\scriptscriptstyle S} \right) $ jest \textbf{homomorfizmem grup addytywnych}. W szczególności  
			\[
			f\left( \mathbf{0}_{\scriptscriptstyle R} = \mathbf{0}_{\scriptscriptstyle S} \right) \text{ oraz } f\left( -r \right) = -f\left( r \right) 
			.\] 
		\item Jeśli  $R$ i $S$ mają $\mathbf{1}$, to niekoniecznie $f\left( \mathbf{1}_{\scriptscriptstyle R} \right) = \mathbf{1}_{\scriptscriptstyle S}$. Na przykład może być $f$ zerowy, tzn $\left(\forall r \in R \right)\left( f\left( r \right) = 0 \right)  $
	\end{enumerate}
\end{remark} 
\textbf{Jako zadanie pokazać:}

Niech $f: R \to S$ będzie homomorfizmem, $R,S$ będą pierścieniami przemiennymi z $\mathbf{1}$ a $S$ będzie dziedziną. Wtedy
\[
f \text{ nie jest zerowy } \implies f\left( \mathbf{1}_{R} \right) = \mathbf{1}_{\scriptscriptstyle S}
.\] 

\textbf{Przykłady} 

Niech $P$ będzie pierścieniem przemiennym oraz $n\in \N_{>0}$.
\begin{enumerate}[label=\arabic*)]
	\item $r_{n}: \mathbb{Z} \to \mathbb{Z} _{n} $ jest epimorfizmem pierścieni.
	\item $\varphi : R \to R\left[ X \right],$ $\varphi \left( r \right) := rX^{0} $ jest monomorfizmem pierścieni.
	
		Utożsamiamy $R$ z $\varphi \left[ R \right] =\{rX^{0}: r\in R\}  $. Wtedy $R \underset{\scriptscriptstyle\text{podp}}{\subseteq} R\left[ R \right] $.

		\begin{remark}[To wymaga wyjaśnienia]
			Mamy
			\[
			a_0X^{0}+ a_1X^{1}+\ldots + a_{n} X^{n} = \underbrace{a_0X^{X^{0}}}_{=a_0}\cdot X^{0}+ \underbrace{a_1X^{0}}_{=a_1}\cdot X^{1}+ \ldots + \ldots + \underbrace{a_{n} X^{0}}_{=a_{n} }\cdot X^{n}
			.\] 
			Z dokładnością do izomorfizmu nad $R$, $R\left[ X \right] $ to pierścień zawierający $R$ oraz $X\in R$ złożony dokładnie z elementów \[
			a_0 + a_1X^{1} + \ldots + a_{n}X^{n}, n\in \N, a_0,a_1,\ldots,a_{n} \in R
			,\] 
			i o tej własności, że
			\[
			a_0 + a_1X^{1} + \ldots + a_{n} X^{n} = b_0 + b_1X^{1} + \ldots + a_{n} X^{n} \iff \left(\forall i\le n \right)\left( a_{i}= b_{i} \right)  
			.\] 
		\end{remark}
	\item \[
	 	\psi : R\left[ X \right] \to \text{Fun}\left( R,R \right), \psi \left( f \right) := \overline{f}
	.\] 
	$\psi $ jest homomorfizmem, nie musi być ,,1-1'' ani ,,na''.
	Na przykład dla $R = \mathbb{Q} $ mamy 
	\[
	\left| \mathbb{Q} \left[ X \right]  \right| = \aleph_{0}, \left| \text{Fun}\left( \mathbb{Q} ,\mathbb{Q}  \right)  \right|= 2^{\aleph_{0}}> \aleph_{0}
	.\] 
\item Niech $Y$ będzie zbiorem oraz $y \in Y$. Definiujemy \textbf{funkcję ewaluacji} 
	\[
	\text{ev}_{y}: \text{Fun}\left( Y, R \right) \to R
	.\] 
	Wtedy $\text{ev}_{y} := f\left( y \right) $ jest epimorfizmem.
\item Mamy tutaj nadużycie notacji/przeciążenie? (Wymaga wyjaśnienia, wykładowca mówił, że w matematyce tak się przyjęło):
	\[
	R\left[ X \right] \overset{\psi }{\longrightarrow} \text{Fun}\left( R,R \right) \overset{\text{ev}_{r}}{\leftarrow }R
	,\] 
	oraz
	\[
		R\left[ X \right] \underset{\text{ev}_{r}}{\leftarrow } R
	.\] 
	Wtedy $f\left( r \right) := \overline{f}\left( r \right) = \text{ev}_{r}\left( f \right) $ jest epimorfizmem. 
\item Mamy:

\begin{figure}[H] % [h!] to sugestia, by umieścić "tutaj"
        \centering
	\adjustbox{scale = 1.3}{
	\begin{tikzpicture}[node distance = 0.7cm and 1.4cm]
  % Tell it where the nodes are
  \node (A) [label={[xshift=4pt]left:$\alpha :$}] {$R$};
  \node (B) [below=of A] {$r$};
  \node (C) [right=of A] {$\text{End}(( R,+ )) $};
  %\node (D) [below=of C] {$\lambda_{r} $};
   \node (D) at (B -| C) {$\lambda_{r}$};
   % Składnia (B -| C) oznacza: weź współrzędną Y z punktu B i współrzędną X z punktu C
  % Tell it what arrows to draw
  \draw[draw=none] (A)-- node[midway] {\small $\rotatein$} (B);
  \draw[|->] (B)--(D);
  \draw[-stealth] (A)--(C);
  \draw[draw=none] (C)-- node [midway] {\small $\rotatein$} (D);
\end{tikzpicture}}
\end{figure}

Gdzie:
\begin{itemize}
	\item $\alpha $ jest homomorfizmem,
	\item $\lambda_{r}\left( x \right) := rx$,
	\item $\alpha \left( r \right) := \lambda_{r}$.
\end{itemize}
Wtedy
\begin{align*}
	\alpha \left( r_1r_2 \right) \left( x \right) &= \lambda_{r_1r_2}\left( x \right) = r_1r_2x = \lambda_{r_1}\left( r_2x \right) = \lambda_{r_1}\left( \lambda_{r_2}\left( x \right)  \right) = \\
		       &= \left( \lambda_{r_1}\circ \lambda_{r_2} \right) \left( x \right) = \left( \alpha \left( r_1 \right) \circ \alpha \left( r_2 \right)  \right) \left( x \right)  
.\end{align*}
Na przykład dla $R = \mathbb{Z} _{n}$ lub $\mathbb{Z} $, $\alpha $ jest izomorfizmem.
\end{enumerate}







